\section{Dataset and Experimental Design}
\label{sec:dataset}

\subsection{Test Dataset Composition}

Our evaluation uses a curated test set of $n=120$ natural language questions spanning three complexity levels: simple, medium, and complex. The dataset composition is presented in Table \ref{tab:dataset_summary}.

\begin{table}[h]
\centering
\caption{Test Dataset Composition by Complexity Level}
\label{tab:dataset_summary}
\small
\begin{tabular}{lcccc}
\toprule
\textbf{Complexity Level} & \textbf{Count} & \textbf{\%} & \textbf{Avg Tables} & \textbf{Avg Time (ms)} \\
\midrule
Simple & 12 & 10.0\% & 1.0 & 32.56 \\
Medium & 64 & 53.3\% & 1.1 & 54.23 \\
Complex & 44 & 36.7\% & 3.1 & 414.70 \\
\midrule
\textbf{Total} & \textbf{120} & \textbf{100.0\%} & \textbf{1.8} & \textbf{184.23} \\
\bottomrule
\end{tabular}

\footnotesize
\textit{Note}: Baseline execution times measured on Oracle Database 26ai with 10 warm-up executions. Average table counts based on table references in FROM and JOIN clauses.
\end{table}


\subsubsection{Complexity Level Definition}

We define complexity levels based on SQL structural characteristics:

\paragraph{Simple (n=12, 10.0\%)}
\begin{itemize}
    \item Single table queries (no JOINs)
    \item Basic filtering with WHERE clauses
    \item Simple aggregations (COUNT, SUM, AVG)
    \item No subqueries or CTEs
    \item Average execution time: 32.56ms (baseline)
\end{itemize}

\paragraph{Medium (n=64, 53.3\%)}
\begin{itemize}
    \item Primarily single-table queries with moderate filters and aggregations
    \item Multiple filter conditions (AND/OR logic)
    \item Aggregations with GROUP BY
    \item Simple subqueries or CASE expressions
    \item Average execution time: 54.23ms (baseline)
\end{itemize}

\paragraph{Complex (n=44, 36.7\%)}
\begin{itemize}
    \item 4+ table JOINs or complex join logic
    \item Nested subqueries and CTEs (Common Table Expressions)
    \item Window functions or analytical queries
    \item Multiple levels of aggregation
    \item Complex business logic (e.g., fiscal year calculations)
    \item Average execution time: 414.70ms (baseline)
\end{itemize}

\subsection{Question Sources}

The 120 test questions were curated from multiple sources including:
\begin{itemize}
    \item \textbf{Existing SQL Datasets}: Questions adapted from WikiSQL and Spider
    \item \textbf{Oracle Documentation Examples}: Questions from Oracle Database tutorial materials
    \item \textbf{Business Domain Knowledge}: Questions from enterprise analytics requirements
    \item \textbf{Stress Testing}: Questions designed to challenge specific aspects (JOINs, aggregations, temporal logic)
\end{itemize}

\subsection{Schema and Data Characteristics}

\begin{table}[h]
\centering
\caption{Evaluation Database Schema Characteristics}
\label{tab:schema_characteristics}
\footnotesize
\begin{tabular}{lcccc}
\toprule
\textbf{Table} & \textbf{Rows} & \textbf{Columns} & \textbf{Primary Indexes} & \textbf{Foreign Keys} \\
\midrule
employees & 10,000 & 8 & emp\_id & dept\_id \\
departments & 50 & 3 & dept\_id & None \\
salaries & 50,000 & 4 & emp\_id, eff\_date & emp\_id \\
orders & 100,000 & 7 & order\_id & customer\_id \\
customers & 5,000 & 5 & cust\_id & None \\
products & 1,000 & 6 & prod\_id & category\_id \\
categories & 50 & 2 & cat\_id & None \\
order\_items & 350,000 & 4 & order\_id, item\_seq & order\_id, prod\_id \\
locations & 100 & 3 & loc\_id & None \\
suppliers & 500 & 4 & supplier\_id & None \\
time\_dim & 3,650 & 3 & date\_key & None \\
exchange\_rates & 9,125 & 3 & currency\_pair, date\_key & None \\
\bottomrule
\end{tabular}

\footnotesize
\textit{Note}: Total schema size: 525,315 rows across 12 tables. Primary indexes use B-tree structures. All tables contain created\_at timestamp and modified\_at LAST\_MODIFIED columns.
\end{table}


The evaluation database includes 12 tables modeling a typical enterprise HR and sales domain:

\subsubsection{Cardinality}
\begin{itemize}
    \item Employees: 10,000 rows
    \item Departments: 50 rows (low cardinality, common JOIN target)
    \item Salaries: 50,000 rows (1:5 ratio with employees for salary history)
    \item Orders: 100,000 rows (high cardinality, range queries)
    \item Customers: 5,000 rows
\end{itemize}

\subsubsection{Data Characteristics}
\begin{itemize}
    \item Nullable columns: 35 nullable columns across schema (important for NULL handling)
    \item Temporal data: Hire dates, salary effective dates, temporal range queries
    \item String data: Mixed case employee names, case-sensitive department codes
    \item Numeric data: Salary ranges from 30,000 to 500,000 with varied distributions
\end{itemize}

\subsection{Baseline Query Sources}

The 120 baseline queries were:

\begin{enumerate}
    \item Written by experienced SQL developers (10+ years experience)
    \item Reviewed by database administrators for correctness and efficiency
    \item Tested for correctness on sample data
    \item Verified to execute without error in test environment
    \item Represent production-quality SQL that defines ground truth
\end{enumerate}

Baseline queries employ:
\begin{itemize}
    \item Proper indexing strategies (14 B-tree indexes on common predicates)
    \item Appropriate JOIN types (INNER vs. LEFT JOINs based on business logic)
    \item Predicate push-down (filters applied at lowest logical level)
    \item Selective column projections (not SELECT *)
\end{itemize}

\subsection{Exclusion Criteria}

We excluded from evaluation:
\begin{itemize}
    \item Queries with DISTINCT (may generate different result format)
    \item UNION/UNION ALL queries (complex semantic equivalence)
    \item DDL/DML queries (CREATE, INSERT, DELETE only SELECT)
    \item Queries with randomness (DBMS\_RANDOM function)
    \item Queries requiring specific execution plans (hints)
\end{itemize}

\subsection{Reproducibility Checklist}

\begin{table}[h]
\centering
\caption{Reproducibility Checklist and Artifact Availability}
\label{tab:reproducibility_checklist}
\footnotesize
\begin{tabular}{lll}
\toprule
\textbf{Artifact} & \textbf{Location} & \textbf{Format} \\
\midrule
\multicolumn{3}{l}{\textit{Test Data and Schema}} \\
Test questions & experiments/test\_questions.json & JSON (120 items) \\
Target schema & Oracle 26ai TPC-H schema (DB instance) & Oracle schema objects \\
\midrule
\multicolumn{3}{l}{\textit{Evaluation Code}} \\
Evaluation framework & experiments/mcp\_evaluation.py & Python 3.10+ script \\
Visualization code & experiments/visualize\_results.py & Python script \\
MCP HTTP server & mcp-server-http.js & Node.js script \\
\midrule
\multicolumn{3}{l}{\textit{Evaluation Results}} \\
Raw results & experiments/results/mcp\_evaluation\_*.json & JSON (timestamped) \\
Comparison results & experiments/results/cmp\_*.json & JSON (timestamped) \\
Visualization outputs & experiments/results/*.png & PNG charts \\
\midrule
\multicolumn{3}{l}{\textit{Configuration Files}} \\
Benchmark config & experiments/benchmark.conf.json & JSON (MCP endpoint, DB connect) \\
Docker setup & docker-compose.yml & Docker Compose \\
\midrule
\multicolumn{3}{l}{\textit{Documentation}} \\
README & README.md & Markdown \\
This paper & research/main.tex & LaTeX \\
\bottomrule
\end{tabular}

\footnotesize
\textit{Note}: Artifacts are versioned in this repository. Docker Compose provides reproducible setup for Oracle DB and supporting services; evaluation outputs are stored as JSON and PNG files under \texttt{experiments/results/}.
\end{table}


All evaluation artifacts are archived for reproducibility:
\begin{itemize}
    \item Test questions in JSON format (test\_questions.json)
    \item Evaluation and analysis code (mcp\_evaluation.py, visualize\_results.py)
    \item Benchmark configuration (benchmark.conf.json)
    \item Container orchestration setup (docker-compose.yml)
    \item Complete result logs in JSON format with timestamps
    \item Generated PNG visualizations for metric reporting
\end{itemize}

\subsection{Threats to Validity}

\subsubsection{Internal Validity}

\begin{itemize}
    \item \textbf{Question Bias}: Questions may not represent real query distribution; mitigated by diverse sources
    \item \textbf{Schema Bias}: Single database schema may over-represent specific patterns; different schemas may show different results
    \item \textbf{MCP Implementation}: Specific Claude API version; results may vary with different LLM backends
\end{itemize}

\subsubsection{External Validity}

\begin{itemize}
    \item \textbf{Database System Generalization}: Results specific to Oracle; PostgreSQL and MySQL may differ
    \item \textbf{Schema Complexity}: Test schema 12 tables; enterprise schemas may have 100+ tables
    \item \textbf{Data Volume}: 100K rows; larger data may expose different failure modes
    \item \textbf{Domain Specificity}: HR/Sales domain; financial or manufacturing domains may differ
\end{itemize}

\subsubsection{Construct Validity}

\begin{itemize}
    \item \textbf{Semantic Equivalence Definition}: Sorted result comparison may miss subtle differences
    \item \textbf{Accuracy Metric}: Treats all queries equally; complex queries may deserve higher weights
    \item \textbf{Latency Measurement}: Single execution; caching effects may differ on repeated execution
\end{itemize}
